\section*{Chapitre \ref{ch:ml:probabilistic-models-and-uncertainty} --- Modèles probabilistes et incertitude}

\begin{solution}{exo:ml:probabilistic-models-and-uncertainty:1}
$y - q = 1~\%$~: perte $0.95\times1~\% = 0.95~\%$. Avec $y = 1~\%$, $y - q = -1~\%$~: perte $(0.95 - 1)(-1~\%) = 0.05~\%$. Cette
asymétrie est ce qui fait du minimiseur le quantile à 95~\%.
\end{solution}

\begin{solution}{exo:ml:probabilistic-models-and-uncertainty:2}
$\lceil251\times0.9\rceil = 226$\up{e} plus petite parmi 250~; $\lceil5\,041\times0.9\rceil = 4\,537$\up{e} parmi 5\,040.
\end{solution}

\begin{solution}{exo:ml:probabilistic-models-and-uncertainty:3}
Les poids $\propto\mu/\sigma^2$~: $5/80^2 : 5/250^2 = 250^2/80^2 = 9.8 : 1$.
\end{solution}

\begin{solution}{exo:ml:probabilistic-models-and-uncertainty:4}
L'intervalle brut a une largeur fixée sur l'année calme de calibration~; lorsque chaque choc est 1.6 fois plus grand, la même largeur
couvre bien moins de rendements (74.0~\%). Le score normalisé divise par l'échelle prédite, qui monte avec la
volatilité EWMA après la rupture, si bien que la largeur suit le nouveau régime, avec un retard (86.9~\%).
\end{solution}

\begin{solution}{exo:ml:probabilistic-models-and-uncertainty:5}
Presque toute la variance prédictive est l'aléa propre du rendement~; davantage de données resserreraient très peu la prévision. La
partie épistémique est plus grande là où le modèle extrapole~: entrées hors de la plage d'entraînement, nouveaux instruments, petits
échantillons.
\end{solution}

\begin{solution}{exo:ml:probabilistic-models-and-uncertainty:6}
L'erreur type d'un ratio de Sharpe annualisé sur trois ans vaut environ $\sqrt{1/3} = 0.58$~; un écart de 0.3 n'est
que du bruit. Sur des échantillons de vingt ans et quatre graines, la même règle fait passer le ratio de Sharpe de 1.72 à 1.91. Tester
les règles de dimensionnement sur de longs échantillons ou de nombreuses graines.
\end{solution}

\begin{solution}{exo:ml:probabilistic-models-and-uncertainty:7}
\texttt{conformal(gamma=0.02)}~: couverture après la rupture 90.0~\% (90.1~\% avec $\gamma = 0.005$). Un $\gamma$ plus grand réagit plus vite
à une suite de dépassements, mais il déplace davantage le seuil d'un jour sur l'autre, élargissant et resserrant la bande sur le bruit~; un
plus petit est plus lisse et plus lent après une rupture.
\end{solution}

\begin{solution}{exo:ml:probabilistic-models-and-uncertainty:8}
$\ell_\alpha(y, q) = (\mathbf 1\{y < q\} - \alpha)(q - y)$. En posant $q = F^{-1}(\alpha)$, $\alpha = F(z)$, $dz =
dF^{-1}(\alpha)$ et en intégrant par parties, $2\int_0^1\ell_\alpha(y, F^{-1}(\alpha))\,d\alpha = 2\int(\mathbf 1\{y < z\} -
F(z))(z - y)\,dF(z) = \int(F(z) - \mathbf 1\{z\ge y\})^2dz$ (Gneiting et Raftery, 2007, donnent l'identité).
\end{solution}

\begin{solution}{pb:ml:probabilistic-models-and-uncertainty:1}
\begin{enumerate}
\item L'espérance d'un score propre est minimisée par la vraie distribution, si bien qu'aucun prévisionniste ne gagne à biaiser sa prévision~; la
log-vraisemblance, la \omterm{def:ml:probabilistic-models-and-uncertainty:crps}{perte pinball} et le CRPS sont propres.
\item CRPS 78.8--79.1 points de base pour tous, couverture 89.4--91.1~\%.
\item CRPS 135.3--136.3, couverture 85.0--87.6~\% (la valeur vraie 91.3~\%).
\item Le score est dominé par la volatilité, que tous prévoient à partir des mêmes variables persistantes~; le
signal de la prévision de moyenne ne vaut que quelques centièmes de la variance.
\item 0.09~\%.
\item La décomposition en variance aléatoire et épistémique, et un intervalle un peu mieux calibré (91.1~\% et 87.6~\% contre
90.6~\% et 87.0~\% pour une seule graine).
\item L'asymétrie et les queues épaisses grâce à deux composantes (ici avec une couverture de 89.7~\% et 86.1~\%).
\item Hors de la plage d'entraînement, pour de nouveaux instruments, et avec peu d'étiquettes.
\item Couverture brute~: 90.7~\% et 74.0~\%~; couverture normalisée~: 90.3~\% et 86.9~\%.
\item 89.9~\% et 90.1~\%~; elle renonce à la garantie à échantillon fini au profit d'un taux d'erreur à long terme sous dérive.
\item Échangeabilité des scores de calibration et de test.
\item D'après la vitesse des dérives attendues et la tolérance à une largeur qui oscille~: petit (0.005) pour les changements lents.
\item Sur trois ans, les ratios de Sharpe ont une erreur type proche de 0.6, supérieure à l'effet.
\item Lecture~: 1.72~; divisé par la variance EWMA, 1.90~; par celle du réseau, 1.91~; par la vraie variance, 1.93. Moyenne vraie 2.94, divisée par la vraie
variance 3.31.
\item La mise à l'échelle repondère les gains espérés du signal~; plus la moyenne est du signal plutôt que du bruit, plus il y a
à repondérer.
\item \emph{Résultat nommé}~: diviser la lecture de dérive par la variance prédite du réseau fait passer le ratio de Sharpe de 1.72
à 1.91 (1.93 avec la vraie variance)~; après le changement de régime, les intervalles conformes scindés couvrent 74.0~\% (brut) ou 86.9~\%
(normalisé) des 90~\% nominaux, la \omterm{def:ml:probabilistic-models-and-uncertainty:conformal}{prédiction conforme} adaptative 90.1~\%.
\item L'écart-type prédictif ou un intervalle, et sa couverture récente.
\item Suivre le taux de réussite (hit rate) de l'intervalle et la distribution des erreurs normalisées avec les détecteurs de dérive du chapitre~12 et les
alarmes du chapitre~27.
\item Queues trop étroites~: un modèle gaussien sous-couvre les extrêmes~; utiliser des quantiles, un mélange ou des intervalles conformes.
\item Une distribution, dont la prévision ponctuelle n'est qu'un résumé.
\end{enumerate}
\end{solution}

\begin{solution}{iq:ml:probabilistic-models-and-uncertainty:1}
Un score dont l'espérance est minimisée en prévoyant la vraie distribution. L'entraînement avec lui rend les distributions honnêtes
optimales~; l'évaluer permet de comparer des prévisions différentes sans récompenser
ni la surconfiance ni la sous-confiance.
\iqlookfor{la définition et l'argument d'incitation.}
\end{solution}

\begin{solution}{iq:ml:probabilistic-models-and-uncertainty:2}
Ajuster des objectifs quantiles (\omterm{def:ml:probabilistic-models-and-uncertainty:crps}{perte pinball}) aux niveaux de l'intervalle, ou un modèle de moyenne plus un modèle de la valeur absolue des résidus,
et les envelopper dans une calibration conforme sur des données hors échantillon (\omterm{def:ml:probabilistic-models-and-uncertainty:quantile}{régression quantile} conformalisée).
\iqlookfor{pertes quantiles et calibration sur des données hors échantillon.}
\end{solution}

\begin{solution}{iq:ml:probabilistic-models-and-uncertainty:3}
Calculer les scores sur un ensemble de calibration~; le seuil est leur $\lceil(n+1)(1-\alpha)\rceil$-ième plus petit~; les intervalles sont les
valeurs dont le score est inférieur à ce seuil~; la couverture vaut au moins $1 - \alpha$ si les scores de calibration et les nouveaux scores sont échangeables. Elle
échoue lorsque les marchés changent de régime (volatilité, liquidité), puisque les scores de calibration ne décrivent plus les nouveaux~;
utiliser des scores normalisés et des mises à jour adaptatives.
\iqlookfor{la construction, la condition d'échangeabilité et le remède.}
\end{solution}

\begin{solution}{iq:ml:probabilistic-models-and-uncertainty:4}
Proportionnellement à la moyenne divisée par la variance (Kelly ou moyenne-variance), éventuellement réduite~; tester sur de longs historiques ou de nombreux
chemins simulés, car les écarts de ratio de Sharpe entre règles sont petits au regard de leurs erreurs types.
\iqlookfor{moyenne sur variance, et un test avec une puissance suffisante.}
\end{solution}

\begin{solution}{iq:ml:probabilistic-models-and-uncertainty:5}
\omterm{def:ml:probabilistic-models-and-uncertainty:aleatoric}{Incertitude aléatoire}~: le bruit de la cible étant donné les entrées, estimé par une variance prédite. \omterm{def:ml:probabilistic-models-and-uncertainty:aleatoric}{Incertitude épistémique}~: incertitude sur le
modèle, estimée par le désaccord d'un ensemble (ou d'une loi a posteriori). La première persiste avec davantage de données~; la seconde diminue.
\iqlookfor{les définitions et un estimateur pour chacune.}
\end{solution}

\begin{solution}{iq:ml:probabilistic-models-and-uncertainty:6}
$\frac{d}{dq}\E\ell_\alpha(Y, q) = G(q) - \alpha$ pour $G$ continue~: négative en dessous du quantile d'ordre $\alpha$, positive au-dessus, si bien que
le minimum est atteint en $G^{-1}(\alpha)$.
\iqlookfor{la dérivée et l'argument de signe.}
\end{solution}
